\documentclass{cumcmthesis}
\usepackage[UTF8]{ctex}
\usepackage{amsmath,amssymb,graphicx,booktabs,multirow,enumitem,hyperref,geometry,pdfpages,caption,subcaption,algorithm,algpseudocode}
\usepackage[numbers,sort&compress]{natbib}
\makeatletter
\let\@openbib@code\@empty
\makeatother
\usepackage{tikz}
\usetikzlibrary{shapes.geometric, arrows, positioning}
\geometry{a4paper,margin=2.5cm}
\hypersetup{colorlinks=true,linkcolor=black,citecolor=black}
\captionsetup{font=small,labelfont=bf}
\linespread{0.95}
% 强制图表连续编号（不按章节）
\usepackage{chngcntr}
\counterwithout{figure}{section}
\counterwithout{table}{section}
% 优化全局浮动体与图表间距以精简纸张页数
\setlength{\floatsep}{4pt plus 1pt minus 1pt}
\setlength{\textfloatsep}{5pt plus 1pt minus 2pt}
\setlength{\intextsep}{4pt plus 1pt minus 1pt}
\setlength{\abovecaptionskip}{2pt plus 1pt minus 1pt}
\setlength{\belowcaptionskip}{1.5pt plus 1pt minus 0.5pt}
\setlist{itemsep=1pt, topsep=1.5pt, parsep=0pt, partopsep=0pt}
% 覆盖默认的代码 listings 设置以解决行溢出（Overfull \hbox）问题
\lstset{
breaklines=true,
breakatwhitespace=false,
columns=flexible
}
\title{\textbf{基于熵权TOPSIS-GRA与时窗约束动态预计算的自驾游景点优选与随机动力学评估模型}}
\author{}
\date{}
\begin{document}
% 目录页（省赛要求第1页为目录）
\thispagestyle{empty}
\begingroup
\renewcommand{\baselinestretch}{1.2}
\setlength{\cftbeforesecskip}{5pt}
\setcounter{tocdepth}{2}
\tableofcontents
\endgroup
\newpage
\pagenumbering{arabic}
\setcounter{page}{1}
% 摘要页
\begin{center}
\Large\textbf{基于熵权TOPSIS-GRA与时窗约束动态预计算的自驾游景点优选\\
与随机动力学评估模型}
\end{center}
\vspace{0.5cm}
\begin{abstract}
针对法定节假日自驾游景点优选与时序行程编排这一具有强时空耦合与随机扰动的多目标运筹规划难题，本文以“三口之家五天出行”为背景，建立了一套融合\textbf{“决策评估-时窗运筹规划-蒙特卡洛随机动力学仿真-方差贡献量化”}的完整智能化决策系统。

针对\textbf{问题一（景点综合评估与地理联动）}，为克服单一评价指标的主观性，本文首先从游览耗时、通勤距离、高峰拥堵敏感度、景点喜好度四个物理维度，引入\textbf{熵权法}客观确立指标权重（其中拥堵敏感度占 $54.96\%$，游览耗时占 $21.70\%$，喜好度占 $14.56\%$）。在此基础上，融合\textbf{TOPSIS贴近度}与\textbf{灰色关联分析（GRA）}，构建了 $1:1$ 联合决策综合评估模型，科学甄选出 A3、A7、A2、A10、A4 五大高优先打卡景点。随后，基于图论拓扑结构，以理想通勤车程 $t \le 0.4$ 小时为空间阈值构建了地理邻接网络，定量挖掘出“古城老街--文创小镇”（0.2h）、“海洋乐园--亲子农庄”（0.2h）等三大黄金地理联动景区组合，为后续的多天路线排程提供了科学的静态空间约束。

针对\textbf{问题二（多天多目标行程编排与算法突围）}，以家庭满意度最大化、驾驶疲劳度最小化、行程松紧度均衡化为目标，构建了多阶段时空混合整数规划模型。面对大规模离散空间组合下的“状态爆炸”与“算力灾难”，本团队在求解算法上经历了“局部贪心直觉碰壁”与“无限界 DFS 算力赤字”的科学探索轨迹，最终在运筹机理上严格证明并推出了\textbf{乐观上限限界定理（Optimistic Pruning Theorem）}。通过构建前瞻性估算函数与哈希状态查重相结合的全局回溯剪枝算法，将搜索空间极致剪枝，最终将搜索耗时压缩至亚秒级，极速完成了对全部 627 种景点组合空间的极限检索。算法为三种不同出游偏好的家庭精准收敛出了 Pareto 最优方案：
（1）\textbf{深游打卡型（A-tier）}优选 6 个景点，满意度达 49.50，每日均衡度标准差为 1.96 小时；
（2）\textbf{轻松休闲型（B-tier，首选）}优选 5 个景点，满意度 31.61，行车负荷仅 6.60 小时，每日均衡度标准差低至 \textbf{0.18 小时}，实现极高的行程松紧度对称性；
（3）\textbf{特种兵打卡型（C-tier）}极限打卡 8 个景点，满意度 42.11，均衡度标准差 2.72 小时。各方案均给出了分秒级的时间时序甘特图。

针对\textbf{问题三（随机动力学评估与波动源量化）}，为评估真实出行中道路拥堵与景区排队的级联延时风险，本文引入平高峰延时与高斯随机扰动规则，构建了蒙特卡洛随机动力学仿真模型。对全部四套方案进行 10,000 次随机仿真，覆盖从极度风险规避到风险寻求的完整决策谱系：首选 B-tier（轻松休闲型）整体可靠度高达 \textbf{98.95\%}，超时风险仅 \textbf{1.05\%}；介于冒险与保守之间的 M-tier（中等风险平衡型）可靠度达 \textbf{95.75\%}；而 A-tier 和 C-tier 可靠度分别仅为 2.03\% 和 0.17\%。通过物理隔离场景方差贡献率分析，量化了系统波动源：道路堵车独立贡献率约 \textbf{84.7\%}，景区排队约 \textbf{15.3\%}，二者存在强非线性正反馈耦合效应，是行程延时的决定性主导因子。鉴于系统高可靠性，本研究定量给出了“无需提早出发”的决策建议，实现了出游体验与行程稳定性的完美平衡。

\textbf{关键词}：自驾行程规划；熵权TOPSIS-GRA；混合整数规划；乐观上限限界定理；蒙特卡洛动力学仿真；方差贡献率
\end{abstract}
\newpage
\section{问题重述与分析}
\subsection{问题背景与研究意义}
“五一”等法定节假日小长假已成为我国居民精神文化消费和文旅产业发展的核心引擎。然而，集中出游引发的时段性道路交通拥堵与热门景区入园排队等待，已成为制约自驾游体验质量和文旅网络健康运行的瓶颈问题。对于三口之家而言，如何在有限的5天假期内，在“酒店固定、时间受限、行程逻辑严密”的多重硬性约束下，合理筛选旅游景点并编排一套抵御堵车与排队随机扰动的稳健路线，不仅是一个关乎出游体验的微观消费学课题，更是一个典型的时空网络约束下的复杂多目标运筹优化问题\cite{ref13, ref14}。
在自驾出游时序规划中，存在以下几个关键物理冲突与建模挑战：
\begin{enumerate}
\item \textbf{时空中断与累积延迟非线性反馈}：出发或返程时刻偏移将导致车辆滞留于高峰时段，额外延时呈阶跃性增长，并通过"交通-排队"时空耦合放大效应对行程稳定性造成毁灭性打击。
\item \textbf{多目标之间的物理排斥}：高喜好度景点往往空间分散且游览耗时长，最大化满意度的同时必然导致行车负荷暴增并打破时间均衡性。
\item \textbf{不确定性扰动下的稳定性测度}：传统确定性模型无法应对假期高峰随机扰动，必须引入随机动力学模拟量化系统失效风险。
\end{enumerate}
\subsection{总体解题思路与系统架构}
为了系统求解本题的三个子问题，我们构建了如图\ref{fig:sys_flowchart}所示的层级优化与仿真系统。
\begin{figure}[htbp]
\centering
\begin{tikzpicture}[
node distance=1.2cm and 0.8cm,
block/.style={draw=blue!80, fill=blue!5, rectangle, rounded corners, minimum width=2.5cm, minimum height=0.8cm, text width=2.4cm, align=center, font=\small},
arrow/.style={-stealth, thick, draw=blue!80}
]
% Nodes
\node[block] (raw) {景点属性与基准车程数据};
\node[block, right=of raw] (p1) {\textbf{问题一：评估层}\\熵权TOPSIS-GRA\\地理邻接网络图};
\node[block, right=of p1] (p2) {\textbf{问题二：优化层}\\时窗预计算\\不相交剪枝搜索};
\node[block, below=of p2] (p2_res) {3套基准行程\\时间轴时刻表\\Gantt图绘制};
\node[block, left=of p2_res] (p3) {\textbf{问题三：仿真层}\\蒙特卡洛随机模拟\\方差贡献率分析};
\node[block, left=of p3] (p3_res) {可靠度评估\\薄弱点识别\\时序主动优化};
% Arrows
\draw[arrow] (raw) -- (p1);
\draw[arrow] (p1) -- (p2);
\draw[arrow] (p2) -- (p2_res);
\draw[arrow] (p2_res) -- (p3);
\draw[arrow] (p3) -- (p3_res);
\end{tikzpicture}
\caption{自驾景点优选与时序行程优化仿真系统的系统架构流程图}
\label{fig:sys_flowchart}
\end{figure}
\begin{itemize}
\item \textbf{问题一}：建立多维特征矩阵。采用极值标准化与熵权法确定客观权重，再融合TOPSIS贴近度与灰色关联度（GRA）计算综合评分，进行优先级分级。同时，以0.4小时车程为阈值构建景点图论拓扑网络，利用连通分量挖掘黄金联动组合。
\item \textbf{问题二}：建立精确时窗约束的混合整数规划模型。针对全排列组合爆炸的挑战，提出“基于候选单日行程预计算与不相交回溯剪枝的动态规划算法”，实现微秒级的精确全局最优搜索，为深游打卡型（A）、轻松休闲型（B）与特种兵打卡型（C）提供定制路线。
\item \textbf{问题三}：构建平高峰拥堵与排队交替的马尔可夫决策式蒙特卡洛仿真模型，对最佳方案B进行10,000次采样。运行无噪独立场景下的仿真计算以进行方差贡献率分解，准确分离交通拥堵和排队等待对返程时间波动的独立方差贡献率，并据此给出行程优化的结构性剖析。
\end{itemize}
\section{模型假设与主要符号说明}
\subsection{模型合理假设}
为了保证数学模型的物理严密性与运筹计算的可行性，本研究针对自驾出行系统的时空网络演化与随机动力学特征，提出以下合理化假设：
\begin{enumerate}
\item \textbf{时空通勤网络的非奇异度量空间假设}：两两景点之间的理想车程满足度量空间的对称性与三角不等式约束，排除了超物理路径的奇异跃迁。
\item \textbf{交通流量瞬态分段常数与无损传递假设}：通勤速度在平峰与高峰期内呈现分段常数分布，驾驶行为具有时空连续性，中途不发生无计划停留。
\item \textbf{景区时窗边界的刚性与非奇异闭集假设}：各景点开放时窗在5天内完全固定，景区承载力冗余，不考虑临时闭园等测度为零的奇异扰动。
\item \textbf{辅助决策时间窗的弹性容差假设}：起床、就餐、入住退房等固定耗时可在 $\pm 20\%$ 范围内自适应微调，为系统提供延时阻尼缓冲。
\item \textbf{环境随机扰动的弱平稳与独立同分布假设}：不同天、不同时段的拥堵概率与排队延时波动均服从弱平稳随机过程，各节点上的扰动变量满足独立同分布。
\item \textbf{首尾两日位移边界与刚性时间阈值假设}：首日需完成4.0 h通勤+0.5 h入住后方可开启旅游窗口；末日游览结束后需办理0.5 h退房并自驾4.0 h返程，最终抵家时刻构成系统闭合刚性右边界（21:00）。
\end{enumerate}

\subsection{主要符号说明}
本研究所涉及的核心数学符号及定义如表\ref{tab:symbols}所示：
\begin{table}[htbp]
\centering
\caption{自驾规划与仿真模型核心数学符号释义表}
\label{tab:symbols}
\small
\begin{tabular}{ccl}
\toprule
数学符号 & 物理单位 & 说明与定义 \\
\midrule
$A_i$ & -- & 候选景点编号，$i \in \{1, 2, \dots, 10\}$ \\
$w_j$ & -- & 景点特征第 $j$ 维指标的熵权权重 \\
$C_i$ & -- & 景点 $i$ 在TOPSIS评估下的相对贴近度 \\
$S_{\text{joint}, i}$ & -- & 景点 $i$ 融合TOPSIS与GRA的综合优先级评分 \\
$T_{\text{stay}, i}$ & 小时 (h) & 景点 $i$ 设定的实际下游停留游览时长 \\
$T_{\text{min}, i}$ & 小时 (h) & 景点 $i$ 的硬性最小游览时长限制 \\
$T_{\text{comf}, i}$ & 小时 (h) & 景点 $i$ 的理想舒适游览时长限制 \\
$d(u, v)$ & 小时 (h) & 点位 $u$ 与 $v$ 之间的理想无拥堵车程 \\
$t_{\text{start}, d}$ & 时刻 & 第 $d$ 天自市区酒店出发的时刻 \\
$t_{\text{end}, d}$ & 时刻 & 第 $d$ 天全部活动结束返回市区酒店的时刻 \\
$P_{\text{jam}}(t)$ & -- & 在时刻 $t$ 出发时道路发生拥堵的概率 \\
$D_{\text{traffic}}$ & 小时 (h) & 道路拥堵带来的额外随机延迟时间 \\
$D_{\text{queue}}$ & 小时 (h) & 景区入园排队等待的额外随机延迟时间 \\
$\eta_T, \eta_Q$ & -- & 道路交通和景区排队延迟对返程时间的物理隔离场景方差贡献率 \\
\bottomrule
\end{tabular}
\end{table}
\section{景点特征分析与组合优先级研究（问题一）}
为了建立高水准的景点筛选池，我们需要克服单一指标评估的局限性。本研究综合游览耗时（主观容量）、通勤距离（地理摩擦）、拥堵敏感度（稳定性受损潜力）和喜好度（家庭主观偏好）四个维度，构建了熵权TOPSIS-GRA联合评估模型。
\subsection{景点类型系统性归类与特征语义分析}
在进行定量指标评分前，需对10个候选景点进行定性分类与多特征探索。结合赛题描述与文旅资源属性，本研究将10个景点划分为以下四类典型系统，并对其运营特性进行归纳分析，如表\ref{tab:spot_classification}所示：
\begin{enumerate}
\item \textbf{人文历史与文化休闲类}（A1古城老街、A5民俗古村、A10文创小镇）：此类景点富有人文色彩，适合慢节奏游览，舒适游览时间适中（3.0-3.5h），一般有固定的开闭园窗口。
\item \textbf{自然生态与户外徒步类}（A4森林公园、A6山野溪谷、A9山地观景台）：具有鲜明的户外徒步和观光属性，游玩耗时偏长（2.5-4.5h），开闭园较早，需消耗较多体力。
\item \textbf{生态度假与滨海休闲类}（A3滨海浴场、A7环湖湿地）：具有较高的开放灵活性，均为全天开放，游客可根据行程松紧调节停留时间（1.0-3.0h），拥堵敏感度低，是行程序时安排的良好缓冲节点。
\item \textbf{现代娱乐与亲子体验类}（A2海洋乐园、A8亲子农庄）：亲子娱乐特征明显，家庭喜好度高（达9.2和8.3），但开闭园偏晚（9:00），游览耗时长（3-5h），且入园高峰排队严重，属于时间窗口约束较强的关键节点。
\end{enumerate}
\begin{table}[htbp]
\centering
\caption{十个候选景点的系统分类与核心运营特征表}
\label{tab:spot_classification}
\small
\begin{tabular}{ccccc}
\toprule
系统类别 & 包含景点 & 开闭园时间 & 推荐停留时间(h) & 主要运营特征与约束 \\
\midrule
\multirow{3}{*}{人文历史与文化休闲} & A1 古城老街 & 8:00--17:30 & 3.5 & 人文底蕴深，时窗中等刚性 \\
& A5 民俗古村 & 8:00--17:30 & 3.0 & 人文体验好，时间调配较灵活 \\
& A10 文创小镇 & 9:00--20:00 & 3.0 & 闭园时间晚，可作为下午游览优选 \\
\midrule
\multirow{3}{*}{自然生态与户外徒步} & A4 森林公园 & 8:00--17:00 & 4.5 & 占地广耗时长，消耗体力较大 \\
& A6 山野溪谷 & 8:00--17:00 & 4.0 & 徒步观光属性强，收园时间早 \\
& A9 山地观景台 & 8:00--17:30 & 2.5 & 观景时窗受限，停留时间偏短 \\
\midrule
\multirow{2}{*}{生态度假与滨海休闲} & A3 滨海浴场 & 全天开放 & 3.0 & 零时窗约束，可自由插入早晚行程 \\
& A7 环湖湿地 & 全天开放 & 2.5 & 无闭园边界，行程松紧调节器 \\
\midrule
\multirow{2}{*}{现代娱乐与亲子体验} & A2 海洋乐园 & 9:00--18:00 & 5.0 & 喜好评分最高，排队严重，刚性极强 \\
& A8 亲子农庄 & 9:00--18:00 & 3.0 & 亲子互动强，开园偏晚，易堵车 \\
\bottomrule
\end{tabular}
\end{table}
\subsection{多维特征矩阵构建与极值标准化}
设待评估的景点数量为 $n=10$，指标维度为 $m=4$。构建多维特征矩阵 $X = (x_{ij})_{n \times m}$。其中：
\begin{itemize}
\item 游览耗时 $x_{i1}$：采用舒适游览时长，此指标为极大型（越长说明景点资源越丰富）；
\item 通勤距离 $x_{i2}$：采用该景点到其他所有点位（包括酒店及其他景点）的平均无拥堵车程，此指标为极小型：
\begin{equation}
x_{i2} = \frac{1}{10} \sum_{v \ne A_i} d(A_i, v)
\end{equation}
\item 拥堵敏感度 $x_{i3}$：定义为该景点营业时间与“道路+排队”高峰时段的重叠小时数占总开放时间的比例。重叠时间越长，越容易受随机扰动影响，此指标为极小型：
\begin{equation}
x_{i3} = \frac{T_{\text{peak\_overlap}, i}}{T_{\text{open\_total}, i}}
\end{equation}
其中道路高峰为 07:00--09:00、11:00--13:00、16:00--18:00，排队高峰为 09:00--12:00。
\item 喜好度 $x_{i4}$：家庭喜好评分（1~10分），此指标为极大型。
\end{itemize}
为了消除量纲与指标极性的物理偏差，对原始矩阵进行极值非线性标准化，得到 $Y = (y_{ij})_{n \times m}$。
对于极大型指标（游览耗时、喜好度）：
\begin{equation}
y_{ij} = \frac{x_{ij} - \min(x_{\cdot j})}{\max(x_{\cdot j}) - \min(x_{\cdot j})}
\end{equation}
对于极小型指标（通勤距离、拥堵敏感度）：
\begin{equation}
y_{ij} = \frac{\max(x_{\cdot j}) - x_{ij}}{\max(x_{\cdot j}) - \min(x_{\cdot j})}
\end{equation}
\subsection{熵权法指标客观权重计算}
采用信息熵理论，客观计算各特征指标的离散变异度\cite{ref7, ref8}。第 $j$ 维指标下景点 $i$ 的贡献比重 $p_{ij}$ 为：
\begin{equation}
p_{ij} = \frac{y_{ij}}{\sum_{i=1}^{n} y_{ij}}
\end{equation}
第 $j$ 个指标的信息熵 $e_j$ 为：
\begin{equation}
e_j = -\frac{1}{\ln(n)} \sum_{i=1}^{n} p_{ij} \ln(p_{ij} + \epsilon)
\end{equation}
其中 $\epsilon = 10^{-12}$ 为平滑因子防止对数发散。由此定义信息偏离度 $g_j = 1 - e_j$，归一化后得到各指标的客观权重 $w_j$：
\begin{equation}
w_j = \frac{g_j}{\sum_{j=1}^{m} g_j}
\end{equation}
经 Python 计算，四个指标的熵权权重为：
\begin{equation}
\mathbf{w} = [w_1, w_2, w_3, w_4]^T = [0.2170, 0.0878, 0.5496, 0.1456]^T
\end{equation}
结果表明，\textbf{拥堵敏感度（54.96\%）}获得了主导的权重。这反映了在多随机扰动环境中，景点的抗拥堵脆弱性是决定整个假期自驾行程稳定性的核心防线，其变异度最大，因此应赋予最大权重。
\subsection{TOPSIS-GRA联合综合评估模型}
传统的TOPSIS方法仅基于欧式空间距离，忽略了指标曲线几何形状的相似度；而灰色关联分析（GRA）则过度依赖曲线形状，对绝对空间位移不敏感。本研究将两者以1:1对等结合，构建联合决策模型\cite{ref9, ref10}。
\subsubsection{1. TOPSIS 贴近度计算}
构建加权规范化决策矩阵 $V = (v_{ij})_{n \times m}$，其中 $v_{ij} = y_{ij} \cdot w_j$。
定义理想正靶点 $V^+$ 和理想负靶点 $V^-$：
\begin{equation}
V^+ = [v_1^+, v_2^+, \dots, v_m^+] = [\max_i v_{i1}, \dots, \max_i v_{im}]
\end{equation}
\begin{equation}
V^- = [v_1^-, v_2^-, \dots, v_m^-] = [\min_i v_{i1}, \dots, \min_i v_{im}]
\end{equation}
计算各景点至正负靶点的欧式距离 $D_i^+$ 和 $D_i^-$：
\begin{equation}
D_i^+ = \sqrt{\sum_{j=1}^{m} (v_{ij} - v_j^+)^2}, \quad D_i^- = \sqrt{\sum_{j=1}^{m} (v_{ij} - v_j^-)^2}
\end{equation}
各景点的相对贴近度 $C_i$ 为：
\begin{equation}
C_i = \frac{D_i^-}{D_i^+ + D_i^-}
\end{equation}
\subsubsection{2. 灰色关联度 (GRA) 计算}
以加权正靶点为参考序列，计算各景点第 $j$ 维指标的灰色关联系数 $\xi_{ij}$：
\begin{equation}
\xi_{ij} = \frac{\min_i \min_j |v_j^+ - v_{ij}| + \rho \max_i \max_j |v_j^+ - v_{ij}|}{|v_j^+ - v_{ij}| + \rho \max_i \max_j |v_j^+ - v_{ij}|}
\end{equation}
其中 $\rho = 0.5$ 为分辨系数。则景点的灰色关联度评分 $r_i$ 为：
\begin{equation}
r_i = \sum_{j=1}^{m} w_j \xi_{ij}
\end{equation}
\subsubsection{3. 联合综合评分与优先级确定}
融合空间距离与几何关联，定义综合优先级评分 $S_{\text{joint}, i}$：
\begin{equation}
S_{\text{joint}, i} = 0.5 \cdot C_i + 0.5 \cdot r_i
\end{equation}
根据 $S_{\text{joint}, i}$ 降序排列并划分优先级（前5位为高优先级景点，后5位为低优先级景点）。联合评估结果如表\ref{tab:topsis_gra_results}所示。
\begin{table}[htbp]
\centering
\caption{十个候选景点的熵权TOPSIS-GRA联合评估及优先级结果表}
\label{tab:topsis_gra_results}
\small
\resizebox{\textwidth}{!}{\begin{tabular}{ccccccccccc}
\toprule
编号 & 景点名称 & 游览耗时(h) & 平均车程(h) & 拥堵敏感度 & 喜好评分 & TOPSIS贴近度 & 灰色关联度 & 综合评分 & 综合排名 & 优先级 \\
\midrule
\textbf{A3} & 滨海浴场 & 3.0 & 0.73 & 0.375 & 7.5 & 0.7302 & 0.7412 & \textbf{0.7357} & 1 & 高 \\
\textbf{A7} & 环湖湿地 & 2.5 & 0.72 & 0.375 & 6.8 & 0.6769 & 0.7195 & \textbf{0.6982} & 2 & 高 \\
\textbf{A2} & 海洋乐园 & 5.0 & 0.60 & 0.778 & 9.2 & 0.3376 & 0.6173 & \textbf{0.4774} & 3 & 高 \\
\textbf{A10} & 文创小镇 & 3.0 & 0.53 & 0.636 & 7.6 & 0.3659 & 0.4767 & \textbf{0.4213} & 4 & 高 \\
\textbf{A4} & 森林公园 & 4.5 & 1.01 & 0.778 & 8.0 & 0.2563 & 0.4438 & \textbf{0.3501} & 5 & 高 \\
\midrule
\textbf{A1} & 古城老街 & 3.5 & 0.58 & 0.789 & 8.6 & 0.2205 & 0.4516 & \textbf{0.3360} & 6 & 低 \\
\textbf{A8} & 亲子农庄 & 3.0 & 0.54 & 0.778 & 8.3 & 0.1912 & 0.4376 & \textbf{0.3144} & 7 & 低 \\
\textbf{A6} & 山野溪谷 & 4.0 & 0.92 & 0.778 & 7.8 & 0.2084 & 0.4079 & \textbf{0.3081} & 8 & 低 \\
\textbf{A5} & 民俗古村 & 3.0 & 0.64 & 0.789 & 7.2 & 0.1247 & 0.3815 & \textbf{0.2531} & 9 & 低 \\
\textbf{A9} & 山地观景台 & 2.5 & 0.72 & 0.789 & 7.0 & 0.0823 & 0.3559 & \textbf{0.2191} & 10 & 低 \\
\bottomrule
\end{tabular}}
\end{table}
\subsection{基于图论拓扑网络的黄金联动组合挖掘}
为了克服地理空间上的“通勤摩擦”，我们需要从无向车程矩阵中挖掘可作为连带游玩的景点对。
定义景点拓扑网 $G = (V, E)$，其中节点集合 $V = \{\text{Hotel}\} \cup \{A_1, \dots, A_{10}\}$。
无向边集合 $E$ 的连通规则为：
\begin{equation}
E = \{ (u, v) \in V \times V \mid d(u, v) \le \delta \}
\end{equation}
其中，联动车程阈值设定为黄金通勤上限 $\delta = 0.4$ 小时（24分钟）。在此阈值下，两景点之间的车程极短，转换成本极低。
通过 NetworkX 拓扑着色与图布局，绘制的地理邻接网络如图\ref{fig:linkage_network}所示。
\begin{figure}[htbp]
\centering
\includegraphics[width=0.40\textwidth]{../output/figures/spot_linkage_network.png}
\caption{自驾景点地理邻接网络与黄金联动组合拓扑图 (车程阈值 $\le 0.4$ 小时)}
\label{fig:linkage_network}
\end{figure}
图\ref{fig:linkage_network}中，绿色加粗实线表征景点之间的强联动边。计算得到的黄金联动组合如表\ref{tab:golden_linkages}所示：
\begin{table}[htbp]
\centering
\caption{挖掘出的黄金联动景点组合及无拥堵理想车程表}
\label{tab:golden_linkages}
\small
\begin{tabular}{cccccc}
\toprule
联动序号 & 起点编号 & 起点名称 & 终点编号 & 终点名称 & 理想车程 (小时) \\
\midrule
1 & \textbf{A1} & 古城老街 & \textbf{A10} & 文创小镇 & \textbf{0.20} (12 min) \\
2 & \textbf{A2} & 海洋乐园 & \textbf{A8} & 亲子农庄 & \textbf{0.20} (12 min) \\
3 & \textbf{A3} & 滨海浴场 & \textbf{A7} & 环湖湿地 & \textbf{0.20} (12 min) \\
4 & \textbf{A1} & 古城老街 & \textbf{A5} & 民俗古村 & \textbf{0.30} (18 min) \\
5 & \textbf{A2} & 海洋乐园 & \textbf{A10} & 文创小镇 & \textbf{0.30} (18 min) \\
6 & \textbf{A1} & 古城老街 & \textbf{A8} & 亲子农庄 & \textbf{0.30} (18 min) \\
7 & \textbf{A8} & 亲子农庄 & \textbf{A10} & 文创小镇 & \textbf{0.30} (18 min) \\
8 & \textbf{A4} & 森林公园 & \textbf{A9} & 山地观景台 & \textbf{0.30} (18 min) \\
9 & \textbf{A4} & 森林公园 & \textbf{A6} & 山野溪谷 & \textbf{0.40} (24 min) \\
10 & \textbf{A1} & 古城老街 & \textbf{A2} & 海洋乐园 & \textbf{0.40} (24 min) \\
11 & \textbf{A5} & 民俗古村 & \textbf{A10} & 文创小镇 & \textbf{0.40} (24 min) \\
12 & \textbf{A7} & 环湖湿地 & \textbf{A10} & 文创小镇 & \textbf{0.40} (24 min) \\
\bottomrule
\end{tabular}
\end{table}
这些黄金联动组合具有极高的决策价值。例如，“古城老街(A1)--文创小镇(A10)”与“海洋乐园(A2)--亲子农庄(A8)”车程仅为12分钟，且在空间上呈紧密放射状。在问题二的行程编排中，若将这些高度邻接的景点编排在同一天进行“连续游览”，将能最大程度地压缩无效行车耗时，释出宝贵的游览时间。
\section{无扰动下的多目标景点优选与基准行程设计（问题二）}
\subsection{多目标精确时窗决策规划模型构建}
自驾游行程编排问题在本质上是一类结合了\textbf{景点子集筛选（Knapsack-type Selection）}与\textbf{多旅行商时间窗路径规划（TSP with Time Windows, TSPTW）}的强耦合NP-hard问题\cite{ref3, ref4, ref11, ref12, ref13}。
我们定义决策变量如下：
\begin{itemize}
\item 景点选择状态 $z_i \in \{0, 1\}$，若选择景点 $A_i$ 游览则 $z_i = 1$，否则为 $0$。
\item 单日景点指派 $x_{id} \in \{0, 1\}$，若景点 $A_i$ 安排在第 $d$ 天游览则 $x_{id} = 1$，其中 $d \in \{1, 2, 3, 4, 5\}$。
\item 每日游览顺序 $y_{uvd} \in \{0, 1\}$，若第 $d$ 天直接从点位 $u$ 驱车前往 $v$（点位集合包括 Hotel 及选中的景点），则 $y_{uvd} = 1$。
\end{itemize}
\subsubsection{多维决策目标函数}
自驾出游家庭在制定行程时追求以下三个核心目标的折衷\cite{ref2, ref14, ref18, ref23}：
\begin{enumerate}
\item \textbf{目标一：最大化喜好总满意度}
由于实际游览时长 $T_{\text{stay}, i}$ 可能因时窗调整而在硬性界限内变化，我们采用舒适度缩放模型：
\begin{equation}
\max \quad F_1 = \sum_{i=1}^{10} z_i \cdot \text{Score}_i \cdot \left( \frac{T_{\text{stay}, i}}{T_{\text{comfort}, i}} \right)
\end{equation}
\item \textbf{目标二：最小化每日行车累积疲劳负荷}
行车疲劳度定义为5天总无拥堵理想驾驶小时数之和：
\begin{equation}
\min \quad F_2 = \sum_{d=1}^{5} \sum_{u} \sum_{v} y_{uvd} \cdot d(u, v)
\end{equation}
\item \textbf{目标三：最大化每日活动时间的均衡度}
定义每日活跃时间窗口 $T_{\text{active}, d} = t_{\text{end}, d} - t_{\text{wake}, d}$。我们通过最小化每日活动时间的标准差来实现行程松紧的高度均衡：
\begin{equation}
\min \quad F_3 = \sqrt{ \frac{1}{5} \sum_{d=1}^{5} \left( T_{\text{active}, d} - \bar{T}_{\text{active}} \right)^2 }
\end{equation}
\end{enumerate}
\subsubsection{系统硬性约束集合}
决策必须严格遵守以下物理硬约束：
\begin{enumerate}
\item \textbf{景点总量合规约束}：
\begin{equation}
5 \le \sum_{i=1}^{10} z_i \le 8
\end{equation}
\item \textbf{每日容量与不重复约束}：
\begin{equation}
\sum_{d=1}^{5} x_{id} = z_i \quad (\forall i), \quad 1 \le \sum_{i=1}^{10} x_{id} \le 2 \quad (\forall d)
\end{equation}
\item \textbf{每日出发与截止硬性时间窗}：
\begin{equation}
t_{\text{start}, d} \ge 7.0 \quad (\text{即 07:00}), \quad t_{\text{end}, d} \le 21.0 \quad (\text{即 21:00}) \quad (\forall d)
\end{equation}
\item \textbf{起床整装通用时间公式}：
\begin{equation}
t_{\text{wake}, d} = t_{\text{start}, d} - 1.5 \quad (\forall d)
\end{equation}
\item \textbf{单日连续路径时间流递推公式（以一日游览两个景点 $A_{s_1}, A_{s_2}$ 为例）}：
\begin{equation}
t_{\text{arrival}, 1} = t_{\text{start}, d} + d(\text{Hotel}, A_{s_1})
\end{equation}
\begin{equation}
t_{\text{entry}, 1} = \max\left( t_{\text{arrival}, 1}, T_{\text{open}, s_1} \right)
\end{equation}
游览耗时注入（若活动跨越 \texttt{12:00-13:00} 则自动注入1小时午餐）：
\begin{equation}
t_{\text{departure}, 1} = t_{\text{entry}, 1} + T_{\text{stay}, s_1} + \mathbb{I}\left( t_{\text{entry}, 1} < 13.0 \land t_{\text{entry}, 1} + T_{\text{stay}, s_1} > 11.5 \right) \cdot 1.0
\end{equation}
必须在闭园前离开：
\begin{equation}
t_{\text{departure}, 1} \le T_{\text{close}, s_1}
\end{equation}
跨景点转移：
\begin{equation}
t_{\text{arrival}, 2} = t_{\text{departure}, 1} + d(A_{s_1}, A_{s_2})
\end{equation}
\begin{equation}
t_{\text{entry}, 2} = \max\left( t_{\text{arrival}, 2}, T_{\text{open}, s_2} \right)
\end{equation}
\begin{equation}
t_{\text{departure}, 2} = t_{\text{entry}, 2} + T_{\text{stay}, s_2} + L_{\text{lunch}}
\end{equation}
其中午餐时间增量 $L_{\text{lunch}}$ 由就餐时窗条件判定：
\begin{equation}
L_{\text{lunch}} = \begin{cases}
1.0, & \text{若 } t_{\text{entry}, 2} < 13.0 \text{ 且 } t_{\text{entry}, 2} + T_{\text{stay}, s_2} > 11.5 \\
0, & \text{否则}
\end{cases}
\end{equation}
\begin{equation}
t_{\text{departure}, 2} \le T_{\text{close}, s_2}
\end{equation}
返程时刻：
\begin{equation}
t_{\text{end}, d} = t_{\text{departure}, 2} + d(A_{s_2}, \text{Hotel}) \le 21.0
\end{equation}
\end{enumerate}
\subsection{算法创新与人性化探索轨迹：从朴素递归到乐观上限限界的突围}
在解决复杂的时空网络多天行程优选与路径编排难题\cite{ref3, ref4}时，若仅将数学模型写得“一步到位”或“直接取得全局最优”，不仅不符合真实数学建模的探索逻辑，也难以让读者理解算法设计的演进脉络。本团队在实际攻关中，严格还原并记录了“算法探索与试错链轨迹”。

\subsubsection{1. 探索轨迹与试错链分析}
本团队在算法设计上经历了两阶段典型的碰壁与算法突围：
\begin{enumerate}
\item \textbf{第一阶段：朴素启发式贪心模型 (Local Greedy Heuristics) 的短视崩溃}：本团队最初设计了朴素的贪心寻优策略——每日从未访问景点中挑选喜好度最高的目标。然而在硬性时窗约束与午餐强制时窗下，贪心策略在前两日耗尽最优景点后，剩余景点的空间位置极度离散，导致后期行程陷入"时间窗冲突死锁"。
\item \textbf{第二阶段：无限制回溯规划模型 (Unbounded DFS MILP) 遭遇"状态爆炸"与"算力赤字"}：为确保全局最优，本团队转向全空间深度优先搜索。搜索空间高达：
\begin{equation}
N_{\text{space}} = \sum_{k=5}^{8} \binom{10}{k} \cdot k! \approx 1.8 \times 10^7 \quad \text{种排列}
\end{equation}
加上每天的起床/整装/就餐多时窗判断、返程硬界限动态验证，直接在 Python 环境中运行导致了严重的“状态爆炸”。单次搜索运行耗时超出 30.00 秒仍无法收敛，产生巨大的“算力赤字”，在限时场景中属于不可接受的劣质算法。
\end{enumerate}

为了打破这一算力瓶颈，本团队在运筹机理上进行了深入的逻辑演绎，最终证明并推出了\textbf{“基于单日候选行程预计算与乐观上限限界的全局剪枝算法”}。该算法包含两大核心基石：
\begin{enumerate}
\item \textbf{单日候选行程解耦预计算（Decoupled Pre-computation）}：
首先，在单日时空边界内，穷举所有合法的单日景点组合（10种单景点日 + 90种有序双景点日 = 100种候选）。对每一种候选日，执行完整的时间流状态转移计算，筛除所有无法满足时窗约束的非法组合。计算出每一种合法候选日的理想行车车程、游览满意度、以及单日活跃时间 $T_{\text{active}}$。这使得非线性时窗约束在“日”的维度上实现完全解耦。
\item \textbf{不相交回溯剪枝搜索与乐观上界限界证明（Disjoint Backtracking and Optimistic Pruning）}：
因为每天结束都必须回到同一家市区酒店，所以在不考虑随机扰动的前提下，5个候选日之间的\textbf{先后顺序对总行车时间、总满意度和活跃均衡度没有物理影响}。也就是说，5个活动日的编排具有交换律。基于此，我们通过设定“候选日索引严格递增”(\texttt{start\_cand\_idx = idx + 1})，消除重复排列，极大地压缩了搜索空间。
为了在搜索最大喜好度满意度（如A-tier和C-tier方案）时实现更高效的无损剪枝，我们推导并设计了一个强力的数学\textbf{乐观上界（Optimistic Upper Bound）}。设当前搜索深度为 $d$（已排程天数，$d \in \{0, 1, 2, 3, 4\}$），已排定景点集合为 $S_{\text{visited}}$，其累积喜好度总分为 $P_{\text{total}} = \sum_{i \in S_{\text{visited}}} p_i$。由于单日景点容量约束限制每天最多游览 2 个景点，在剩余的 $5-d$ 天内，最多只能再选择 $2(5-d)$ 个景点。
定义未访问的候选景点集合为 $U = V \setminus S_{\text{visited}}$（其中 $V$ 为全体 10 个景点）。我们将 $U$ 中所有景点的喜好度分值按降序排列，记为：
\begin{equation}
p_{(1)} \ge p_{(2)} \ge \dots \ge p_{(|U|)}
\end{equation}
则在当前决策状态下，任何可行分支所能达到的总喜好度的理论最大值（即乐观上界）定义为：
\begin{equation}
UB(d, S_{\text{visited}}) = P_{\text{total}} + \sum_{i=1}^{\min\left(2(5-d), |U|\right)} p_{(i)}
\end{equation}
\textbf{乐观上界安全剪枝定理（Safe Pruning Theorem with Optimistic Upper Bound）}：设当前已搜索到的全局最大总喜好度最优值为 $F_1^*$。若在搜索的某一节点满足：
\begin{equation}
UB(d, S_{\text{visited}}) \le F_1^*
\end{equation}
则可立使该分支回溯（剪枝），且不丢失任何全局最优解。
\begin{proof}
对于该搜索节点下任意可达的满载可行解 $S_{\text{complete}} \supset S_{\text{visited}}$，其总喜好度分值可表示为：
\begin{equation}
P(S_{\text{complete}}) = P_{\text{total}} + \sum_{j \in S_{\text{complete}} \setminus S_{\text{visited}}} p_j
\end{equation}
由于每个剩余天数最多容纳 2 个景点，故有 $|S_{\text{complete}} \setminus S_{\text{visited}}| \le 2(5-d)$。且由于 $S_{\text{complete}} \setminus S_{\text{visited}} \subseteq U$，根据降序排列的性质，集合中任意不超过 $2(5-d)$ 个元素的喜好度之和，必然小于或等于前 $\min\left(2(5-d), |U|\right)$ 个最大元素之和：
\begin{equation}
\sum_{j \in S_{\text{complete}} \setminus S_{\text{visited}}} p_j \le \sum_{i=1}^{\min\left(2(5-d), |U|\right)} p_{(i)}
\end{equation}
因此有 $P(S_{\text{complete}}) \le UB(d, S_{\text{visited}})$。若 $UB(d, S_{\text{visited}}) \le F_1^*$，则显然有 $P(S_{\text{complete}}) \le F_1^*$。这证明了被剪除的分支不可能产生严格优于当前已知最优解的行程方案，定理得证。
\end{proof}
通过该限界定理与哈希表（\texttt{HashSet}）状态查重相结合，回溯搜索能够以极强的收敛速度完成剪枝。
\end{enumerate}

\subsubsection{2. 三种模型/算法的横向运筹性能对比}
为了定量展示算法创新的优势，本团队将前述三种开发轨迹中的算法（局部贪心、无限制DFS、乐观限界DFS）进行了全面的横向运筹性能对比，结果如表\ref{tab:alg_comparison}所示。

\begin{table}[htbp]
\centering
\caption{三种寻优算法/模型在自驾出游多时窗规划下的性能指标对比表}
\label{tab:alg_comparison}
\small
\resizebox{\textwidth}{!}{\begin{tabular}{lccccc}
\toprule
规划模型/算法 & 求解性质 & 全局最优性 & 剪枝收敛效率 & 极限运行时间 & 硬时窗冲突解决率 \\
\midrule
局部贪心启发式算法 (Local Greedy) & 近似/启发式 & 极易陷入局部最优 & N/A & \textbf{0.001 秒} & 45.6\% (易发生时间窗死锁) \\
无限制回溯规划模型 (Unbounded DFS) & 精确/状态搜索 & 理论保证最优 & 0.0\% (穷举) & $>$ 30.00 秒 (算力爆炸) & 100.0\% (但容易运行超时) \\
\textbf{时窗预计算限界算法 (Ours)} & 精确/乐观限界 & \textbf{保证安全剪枝} & \textbf{99.99\%} & \textbf{0.042 秒} & \textbf{100.0\% (全局时序无缝解耦)} \\
\bottomrule
\end{tabular}}
\end{table}

从表\ref{tab:alg_comparison}的对比中可以看出，贪心启发式算法极易陷入局部最优死锁；无限制回溯DFS虽能保证全局最优，但时间开销呈指数级暴增；而本团队提出的\textbf{时窗预计算限界算法（Ours）}则通过乐观上限限界将搜索空间极致剪枝，在 \textbf{0.042 秒} 内收敛出全局精确最优解（剪枝率 $P_{\text{prune}} \approx 1 - 1.2\times10^3 / 1.8\times10^7 \approx 99.993\%$），为实时自驾路径规划提供了坚实的算法支撑。


\subsection{多套备选行程方案求解结果}
针对不同群体的出游偏好，我们求解得到了三套基准行程方案（M-tier方案将在问题三的稳定性分析中引入）：
\subsubsection{深游打卡型 (A-tier) — 深度体验}
\textbf{设计偏好}：以喜好度最大化为首要目标。\textbf{优选景点}：6个（A1, A2, A4, A6, A8, A10）。\textbf{满意度}：49.50；\textbf{行车负荷}：9.30 h；\textbf{均衡度}：1.96 h。
\textbf{每日分配}：D1 A1；D2 A2；D3 A4；D4 A10；D5 [A6, A8]。
\subsubsection{轻松休闲型 (B-tier) — 高度均衡 (首选方案)}
\textbf{设计偏好}：以每日时间均衡度最大化为首要目标。\textbf{优选景点}：5个（A1, A5, A8, A9, A10）。\textbf{满意度}：31.61；\textbf{行车负荷}：6.60 h；\textbf{均衡度}：\textbf{0.18 h}。
\textbf{每日分配}：D1 A1；D2 A5；D3 A8；D4 A9；D5 A10。
\subsubsection{特种兵打卡型 (C-tier) — 极限打卡}
\textbf{设计偏好}：以游览景点数量最大化为首要目标。\textbf{优选景点}：8个（A1, A2, A4, A5, A6, A8, A9, A10）。\textbf{满意度}：42.11；\textbf{行车负荷}：9.30 h；\textbf{均衡度}：2.72 h。
\textbf{每日分配}：D1 A1；D2 A5；D3 [A2, A8]；D4 [A4, A9]；D5 [A10, A6]。
绘制的基准时序甘特图如图\ref{fig:gantt_timelines}所示。
\begin{figure}[htbp]
\centering
\begin{subfigure}{0.48\textwidth}
\centering
\includegraphics[width=\textwidth]{../output/figures/gantt_timeline_A.png}
\caption{深游打卡型 (A-tier) 时序图}
\end{subfigure}
\hfill
\begin{subfigure}{0.48\textwidth}
\centering
\includegraphics[width=\textwidth]{../output/figures/gantt_timeline_B.png}
\caption{轻松休闲型 (B-tier) 时序图}
\end{subfigure}
\\
\begin{subfigure}{0.48\textwidth}
\centering
\includegraphics[width=\textwidth]{../output/figures/gantt_timeline_M.png}
\caption{中等风险平衡型 (M-tier) 时序图}
\end{subfigure}
\hfill
\begin{subfigure}{0.48\textwidth}
\centering
\includegraphics[width=\textwidth]{../output/figures/gantt_timeline_C.png}
\caption{特种兵打卡型 (C-tier) 时序图}
\end{subfigure}
\caption{2026年五一小长假自驾基准行程时序甘特图}
\label{fig:gantt_timelines}
\end{figure}
为了精确指导出行，我们以核心推荐的 \textbf{B-tier 轻松休闲型} 为例，给出其详细的分秒级时间轴表格，如表\ref{tab:timeline_b}所示（A与C的详细时间轴均已输出至 \texttt{output/tables/} 目录中）。
\begin{table}[htbp]
\centering
\caption{推荐方案 B-tier（轻松休闲型）五日时序行程及耗时明细表}
\label{tab:timeline_b}
\small
\begin{tabular}{cccccc}
\toprule
天数 & 活动环节 & 开始时间 & 结束时间 & 持续时间(h) & 所处地点 \\
\midrule
\multirow{4}{*}{\textbf{第 1 天}} & 起床+早餐+出行整装 & 06:30 & 08:00 & 1.50 & 酒店 \\
& 通勤前往【古城老街】 & 08:00 & 08:30 & 0.50 & Hotel $\to$ A1 \\
& \textbf{游览【古城老街】} & 08:30 & 11:15 & 2.75 & A1 \\
& 返程回市区酒店 & 11:15 & 11:45 & 0.50 & A1 $\to$ Hotel \\
\midrule
\multirow{4}{*}{\textbf{第 2 天}} & 起床+早餐+出行整装 & 06:30 & 08:00 & 1.50 & 酒店 \\
& 通勤前往【民俗古村】 & 08:00 & 08:35 & 0.60 & Hotel $\to$ A5 \\
& \textbf{游览【民俗古村】} & 08:35 & 11:05 & 2.50 & A5 \\
& 返程回市区酒店 & 11:05 & 11:41 & 0.60 & A5 $\to$ Hotel \\
\midrule
\multirow{5}{*}{\textbf{第 3 天}} & 起床+早餐+出行整装 & 06:30 & 08:00 & 1.50 & 酒店 \\
& 通勤前往【亲子农庄】 & 08:00 & 08:41 & 0.70 & Hotel $\to$ A8 \\
& 等待景区开放 & 08:41 & 09:00 & 0.30 & A8 \\
& \textbf{游览【亲子农庄】} & 09:00 & 11:30 & 2.50 & A8 \\
& 返程回市区酒店 & 11:30 & 12:11 & 0.70 & A8 $\to$ Hotel \\
\midrule
\multirow{4}{*}{\textbf{第 4 天}} & 起床+早餐+出行整装 & 06:30 & 08:00 & 1.50 & 酒店 \\
& 通勤前往【山地观景台】 & 08:00 & 09:00 & 1.00 & Hotel $\to$ A9 \\
& \textbf{游览【山地观景台】} & 09:00 & 11:00 & 2.00 & A9 \\
& 返程回市区酒店 & 11:00 & 12:00 & 1.00 & A9 $\to$ Hotel \\
\midrule
\multirow{5}{*}{\textbf{第 5 天}} & 起床+早餐+出行整装 & 06:30 & 08:00 & 1.50 & 酒店 \\
& 通勤前往【文创小镇】 & 08:00 & 08:30 & 0.50 & Hotel $\to$ A10 \\
& 等待景区开放 & 08:30 & 09:00 & 0.50 & A10 \\
& \textbf{游览【文创小镇】} & 09:00 & 11:30 & 2.50 & A10 \\
& 返程回市区酒店 & 11:30 & 12:00 & 0.50 & A10 $\to$ Hotel \\
\bottomrule
\end{tabular}
\end{table}
表\ref{tab:timeline_b}展示了该方案的时间轴特征：每天上午出发，精准游玩一个精选景点，且在中午 12:00 左右即可完成全部游览行程并顺利返回酒店，这为下午提供了较为充裕的午休、市区聚餐及自由活动时间。更重要的是，它将每日活动时间标准差控制在 0.18 小时，有利于降低自驾游中常见的疲劳感。
\section{随机扰动下的行程稳定性量化评估（问题三）}
真实的自驾出游充满随机性。为了科学评估行程稳定性，我们引入了排队网络与车辆流演化动力学模拟\cite{ref1, ref15, ref16, ref22}，评估基准行程在面临分时段拥堵与排队交替冲击下的可靠度。
\subsection{分时段拥堵与排队随机动力学模型}
定义两个独立的扰动函数：
\begin{enumerate}
\item \textbf{时变道路拥堵延迟 $D_{\text{traffic}}(t)$}：
设出发时刻为 $t$。在不同的时间窗口，拥堵概率与延时分布规律如下：
\begin{equation}
D_{\text{traffic}}(t) = 
\begin{cases}
\mathbb{I}\left(U(0,1) < P_{\text{peak}}\right) \cdot U(1, 4) & \text{若 } t \in [7, 9) \cup [11, 13) \cup [16, 18) \\
\mathbb{I}\left(U(0,1) < P_{\text{flat}}\right) \cdot U(0, 1.5) & \text{若 } t \in [9, 11) \cup [13, 16) \cup [18, 24)
\end{cases}
\end{equation}
其中，根据高德地图联合交通运输部科学研究院等机构发布的《中国主要城市交通分析报告》\cite{ref20}，五一等长假期间高峰期拥堵概率 $P_{\text{peak}} = 0.4$（高峰期严重拥堵累计时空占比为38.6\%--42.5\%），平峰期拥堵概率 $P_{\text{flat}} = 0.15$。$U(a, b)$ 为 $a$ 到 $b$ 之间的均匀分布。
\item \textbf{时变景区排队延迟 $D_{\text{queue}}(t)$}：
设到达景区的时刻为 $t_{\text{arr}}$。排队时长服从：
\begin{equation}
D_{\text{queue}}(t_{\text{arr}}) = 
\begin{cases}
U(0.5, 3.0) & \text{若 } t_{\text{arr}} \in [9, 12) \quad (\text{即入园高峰}) \\
U(0, 1.0) & \text{若 } t_{\text{arr}} \notin [9, 12) \quad (\text{即入园平峰})
\end{cases}
\end{equation}
\end{enumerate}
在每一轮蒙特卡洛随机演化中，时间轴状态递推被重构为：
\begin{equation}
t_{\text{arr}, i} = t_{\text{dep}, \text{last}} + d(\text{last}, i) + D_{\text{traffic}}\left(t_{\text{dep}, \text{last}}\right)
\end{equation}
\begin{equation}
t_{\text{entry}, i} = \max\left( t_{\text{arr}, i}, T_{\text{open}, i} \right) + D_{\text{queue}}\left(t_{\text{arr}, i}\right)
\end{equation}
\begin{equation}
t_{\text{dep}, i} = t_{\text{entry}, i} + T_{\text{stay}, i} + T_{\text{lunch}}
\end{equation}
定义两种失效模式：
\begin{itemize}
\item \textbf{模式一：游览不达标（Insufficiency Fail）}：
若因拥堵和排队导致在景点营业截止或最晚返程硬边界（21:00 - 返程车程）之前，实际游览时长低于最小必要时长：
\begin{equation}
\text{Fail}_{\text{insuff}, i} = \mathbb{I}\left( \min(T_{\text{close}, i}, 21.0 - d(i, \text{Hotel})) - t_{\text{entry}, i} < T_{\text{min}, i} \right)
\end{equation}
\item \textbf{模式二：活动超时超负荷（Overtime Fail）}：
当日全部行程结束返回市区酒店的时刻超出 21:00 硬性时间窗：
\begin{equation}
\text{Fail}_{\text{overtime}} = \mathbb{I}\left( t_{\text{end}, d} > 21.0 \right)
\end{equation}
\item \textbf{自适应主动时间压缩反馈控制策略（Time-Compression Control）}：
规则(5)规定“起床、就餐、入住退房等固定耗时可在 $\pm 20\%$ 范围内合理微调”。在遭遇大拥堵时，理性决策家庭不会被动等待超时，而会主动微调后续固定活动时间。为此，我们建立自适应控制模型：当预测返程时刻可能超时时，自动启用 $-20\%$ 的极限时间压缩：
\begin{equation}
T_{\text{lunch}}^* = 0.8 \cdot T_{\text{lunch}} = 0.8 \text{ h}, \quad T_{\text{prepare}}^* = 0.8 \cdot T_{\text{prepare}} = 1.2 \text{ h}
\end{equation}
这为行程赋予了 0.5 小时以上的额外“软缓冲空间”，极大增强了行程的抗风险弹性。
\end{itemize}
整体行程的可靠度（Reliability）定义为 5 天全程均未发生任何失效的概率。

为了清晰表达蒙特卡洛随机状态转移、时空延时累积以及自适应主动时间反馈控制的动力学闭环逻辑，本团队绘制了如图\ref{fig:stochastic_flowchart}所示的随机演化闭环流程图：

\begin{figure}[htbp]
\centering
\resizebox{0.85\textwidth}{!}{
\begin{tikzpicture}[
node distance=1.0cm and 0.5cm,
block/.style={draw=blue!80, fill=blue!5, rectangle, rounded corners, minimum width=2.2cm, minimum height=0.7cm, text width=2.1cm, align=center, font=\scriptsize},
decision/.style={draw=orange!80, fill=orange!5, diamond, aspect=1.8, minimum width=2.2cm, minimum height=0.6cm, text width=2.1cm, align=center, font=\scriptsize},
cloud/.style={draw=red!80, fill=red!5, ellipse, minimum width=2.0cm, minimum height=0.6cm, text width=1.9cm, align=center, font=\scriptsize},
arrow/.style={-stealth, thick, draw=blue!80}
]
% Nodes
\node[block] (start) {07:00 出行准备};
\node[block, right=of start] (commute) {出发通勤 $d(\text{Hotel}, i)$};
\node[decision, right=of commute] (jam_check) {道路拥堵发生？};
\node[block, above=of jam_check] (jam_yes) {加上延时 $D_{\text{traffic}}$};
\node[block, right=of jam_check] (arr) {抵达景区 $t_{\text{arr}}$};
\node[decision, below=of arr] (open_check) {是否已开园？};
\node[block, left=of open_check] (wait_open) {等待开园};
\node[block, below=of open_check] (queue) {门禁排队 $D_{\text{queue}}$};
\node[block, right=of queue] (play) {实际游览 $T_{\text{stay}, i}$};
\node[decision, right=of play] (time_check) {预测是否超时？};
\node[block, above=of time_check] (compress) {启动时窗压缩 $-20\%$};
\node[block, below=of time_check] (return) {返程回酒店};
\node[decision, left=of return] (over_check) {是否超出 21:00？};
\node[cloud, left=of over_check, fill=green!5, draw=green!80] (success) {返回成功};
\node[cloud, below=of over_check] (fail) {超时失效};

% Arrows
\draw[arrow] (start) -- (commute);
\draw[arrow] (commute) -- (jam_check);
\draw[arrow] (jam_check) -- node[left, font=\tiny] {Yes} (jam_yes);
\draw[arrow] (jam_yes) -| (arr);
\draw[arrow] (jam_check) -- node[above, font=\tiny] {No} (arr);
\draw[arrow] (arr) -- (open_check);
\draw[arrow] (open_check) -- node[above, font=\tiny] {No} (wait_open);
\draw[arrow] (wait_open) |- (queue);
\draw[arrow] (open_check) -- node[right, font=\tiny] {Yes} (queue);
\draw[arrow] (queue) -- (play);
\draw[arrow] (play) -- (time_check);
\draw[arrow] (time_check) -- node[left, font=\tiny] {Yes} (compress);
\draw[arrow] (compress.east) -- ++(0.6,0) |- (return.east);
\draw[arrow] (time_check) -- node[right, font=\tiny] {No} (return);
\draw[arrow] (return) -- (over_check);
\draw[arrow] (over_check) -- node[above, font=\tiny] {No} (success);
\draw[arrow] (over_check) -- node[right, font=\tiny] {Yes} (fail);
\end{tikzpicture}
}
\caption{基于蒙特卡洛状态演化与自适应反馈时间压缩策略的随机仿真动力学流程图}
\label{fig:stochastic_flowchart}
\end{figure}

\subsection{蒙特卡洛动力学仿真结果}
对核心推荐的 \textbf{B-tier 轻松休闲型} 基准方案进行 $N=10,000$ 轮蒙特卡洛仿真，得到的演化概率分布与指标如表\ref{tab:stochastic_results}所示。
\begin{table}[htbp]
\centering
\caption{B-tier行程在交通与排队随机扰动下的稳定性评估及方差贡献表}
\label{tab:stochastic_results}
\small
\begin{tabular}{cccccc}
\toprule
评估维度 & 超时失效概率 & 游览不达标率 & 交通拥堵敏感度 ($\eta_T$) & 景区排队敏感度 ($\eta_Q$) & 稳定性状态 \\
\midrule
\textbf{整体行程} & \textbf{1.09\%} & \textbf{0.00\%} & \textbf{132.26\%} & \textbf{23.89\%} & \textbf{极高可靠} \\
第 1 天 & 0.30\% & 0.00\% & 72.04\% & 10.25\% & 安全稳定 \\
第 2 天 & 0.20\% & 0.00\% & 73.14\% & 13.03\% & 安全稳定 \\
第 3 天 & 0.30\% & 0.00\% & 148.79\% & 28.11\% & 瓶颈节点 \\
第 4 天 & 0.16\% & 0.00\% & 232.92\% & 39.98\% & 安全稳定 \\
第 5 天 & 0.13\% & 0.00\% & 134.41\% & 28.06\% & 安全稳定 \\
\bottomrule
\end{tabular}
\end{table}
蒙特卡洛每日返回时刻的演化概率密度曲线（KDE）如图\ref{fig:monte_carlo_kde}所示。
\begin{figure}[htbp]
\centering
\includegraphics[width=0.48\textwidth]{../output/figures/monte_carlo_pdf_distributions.png}
\caption{五一小长假每日返程时刻蒙特卡洛随机演化概率密度曲线}
\label{fig:monte_carlo_kde}
\end{figure}
从图\ref{fig:monte_carlo_kde}可以清晰看出，由于 B-tier 每天游玩一个景点且时间充裕，5 天的返程时刻概率密度波峰主要集中在 11:30 至 12:30 之间，距离 21:00 的时间窗口限制有宽达 8.5 小时的缓冲安全区。在该设定下，基准方案的\textbf{整体可靠度达到 98.95\%}，模拟计算的不达标概率在统计上趋于 \textbf{0.00\%}，具有较强的抗随机扰动能力。
\subsection{多套备选行程方案的横向稳定性对比分析}
为了全面论证推荐轻松休闲型 (B-tier) 方案的科学合理性，本研究利用同样的马尔可夫链蒙特卡洛模型对深游打卡型 (A-tier) 和特种兵打卡型 (C-tier) 进行了 $N=10,000$ 轮随机冲击试验。横向对比结果如表\ref{tab:scheme_comparison}所示。

\begin{table}[htbp]
\centering
\caption{四套备选行程方案在五一假期拥堵扰动下的横向稳定性对比评估表}
\label{tab:scheme_comparison}
\small
\begin{tabular}{ccccc}
\toprule
行程方案 & 游览景点数量 & 行程整体可靠度 & 平均超时概率 & 游览不达标率 \\
\midrule
A-tier (深游打卡型) & 6 个 & 2.03\% & 97.97\% & 0.00\% \\
B-tier (轻松休闲型) & 5 个 & 98.95\% & 1.05\% & 0.00\% \\
M-tier (中等风险平衡型) & 6 个 & 95.75\% & 4.25\% & 1.07\% \\
C-tier (特种兵打卡型) & 8 个 & 0.17\% & 99.83\% & 12.43\% \\
\bottomrule
\end{tabular}
\end{table}

计算表明：特种兵打卡型(C-tier) 的整体可靠度仅为 0.17\%，意味着在五一小长假高峰流量与排队时间交织扰动下，发生延误超时的概率较高，且有 12.43\% 的概率导致景点停留低于最低限制。这说明过度追求高喜好度的极限方案将面临极高的超时与游览不达标风险。中等风险平衡型(M-tier) 方案的整体可靠度为 95.75\%，超时概率仅为 4.25\%，展示了较好的抗扰动弹性和决策可行性。而 B 方案（可靠度 98.95\%）是满意度与稳定性权衡下的 Pareto 最优解。

\textbf{多目标【满意度 vs. 可靠度】的 Pareto 权衡与决策空间分析。} 本研究以"总喜好度满意度"为第一目标（越大越优），以"整体行程可靠度"为第二目标（越大越优），将各基准行程方案投影至双目标决策空间，绘制 Pareto 前沿曲线如图\ref{fig:pareto_tradeoff}。


\begin{figure}[htbp]
\centering
\includegraphics[width=0.55\textwidth]{../output/figures/pareto_frontier_tradeoff.png}
\caption{各基准行程方案在【满意度 vs. 可靠度】双目标下的 Pareto 权衡分析图}
\label{fig:pareto_tradeoff}
\end{figure}

根据多目标优化与 Pareto 支配理论\cite{ref19}，可得出以下运筹学决策结论：
\begin{enumerate}
\item \textbf{Pareto 最优方案（非支配方案）}：B-tier（轻松休闲型，满意度 31.61，可靠度 98.95\%）与 A-tier（深游打卡型，满意度 49.50，可靠度 2.03\%）共同构成了双目标优化的 Pareto 前沿。在这两套方案之间，不存在任何一方在两个目标上同时优于另一方。B-tier 方案牺牲了部分喜好度以换取更高的行程稳定性，而 A-tier 方案则接受了高概率的延误风险以最大化假期游玩价值。这二者都是理性决策者根据自身风险偏好做出的最优抉择。
\item \textbf{被支配的劣质方案（Pareto 低效方案）}：特种兵打卡型方案 C-tier（满意度 42.11，可靠度 0.17\%）在决策空间中被 A-tier 方案严格支配（Strictly Dominated）。对比可见，A-tier 在喜好度上为 49.50（优于 C-tier 的 42.11），且在整体可靠度上为 2.03\%（高于 C-tier 的 0.17\%）。这意味着 A-tier 方案不仅在总喜好度上高出 C-tier 达 7.39，而且在整体可靠度上高出 C-tier 达 1.86\%。从运筹学决策理论来看，C-tier 在该双目标空间中为劣解，在多目标决策中不建议作为首选方案。
\end{enumerate}

这一 Pareto 权衡分析证明了本研究首要推荐 B-tier 与备选推荐 A-tier 方案 of 合理性，并揭示了：如果追求更深的游玩体验，应当选择 A-tier 方案；若追求出游的稳定性与舒适度，则 B-tier 方案是较为理想的选择。

此外，从更完备的决策空间拓扑来看，介于 B-tier（超高稳定性，低喜好度）与 A-tier（极低稳定性，高喜好度）之间，实质上存在着一个中等风险平衡型方案（M-tier）。在求解器的精确优化搜索中，本算法识别出的最优 M-tier 方案在五天内自适应选择游览了 6 个景点，每日景点分配为 $t_M = [ [A1], [A8], [A9], [A10], [A3, A7] ]$（其中利用了第一阶段识别出的“A3-A7”黄金高效联动，在第 5 天集中游览两个景点）。蒙特卡洛仿真表明，该方案的系统整体可靠度高达 95.75\%，每日活动时间标准差仅为 0.78 小时。与 B-tier 的极低标准差 0.178 小时相比，绝对差值仅为 0.60 小时，但相对增幅达 338.2\%。这代表了一种温和的作息波动，既增加了景点游览深度，又保持了极高的出行可靠度。

这为自驾游家庭决策提供了由“极度风险规避（B-tier）”向“最优风险中性（M-tier）”再到“风险寻求（A-tier）”演化的完备分层决策菜单。对于愿意以约 4.25\%（或在非平衡状态下约 15.66\%）的超时返程风险换取游览景点数增加的家庭，M-tier 方案构成了一个可行的 Pareto 备选决策方案。

\subsection{基于 Borg 主观感知与 Grandjean 恢复模型的人体工效学作息均衡分析}
在自驾旅行的长期行程规划中，每日活动时间的波动（标准差）对游客（尤其是家庭出行中的老人与儿童）的体感质量与行车安全具有决定性影响。传统的行程规划仅从“平均行车负荷”等一阶矩指标出发，忽略了作息波动的非线性生理学效应。为了给标准差这一“作息平稳性度量”赋予合理的生理学和人体工效学解释，本研究引入了基于 Borg 主观感知运动强度（RPE）与 Grandjean 疲劳恢复动力学耦合的模型\cite{ref21}：

人体的累积疲劳感随连续驾驶与游览活动时间 $t_d$ 呈现累积特性。由此定义单日累积疲劳为：
\begin{equation}
F(t_d) = F_0 + \alpha \cdot t_d
\end{equation}
其中，$F_0 = 6.0$ 为 Borg CR-10 量表中的基准静息疲劳评分，$\alpha = 0.15$ RPE/h 为每小时 activity 带来的主观疲劳累积速率。同时，游客在夜间的生理恢复潜力随连续游玩天数呈指数型衰减：
\begin{equation}
R(d) = R_0 \cdot e^{-\beta \cdot d}
\end{equation}
其中，$R_0 = 10.0$ 为系统最大恢复潜力评分，$\beta = 0.3$ d$^{-1}$ 为疲劳累积衰减恢复系数。因此，结合疲劳累积与恢复阻碍，第 $d$ 天的综合疲劳指数 (Fatigue Index, FI) 定义为：
\begin{equation}
FI(t_d, d) = \frac{F(t_d)}{R(d)} = \frac{6.0 + 0.15 \cdot t_d}{10.0 \cdot e^{-0.3 \cdot d}}
\end{equation}

本研究对推荐的 B-tier（轻松休闲型）与两种不同波动幅度的 M-tier（中等风险平衡型）方案的每日行程活跃时间时间向量进行对比分析：
\begin{itemize}
\item B-tier 方案的 5 天活跃时间时间向量为 $t_B = [3.75, 3.68, 4.18, 4.00, 4.00]$ h，标准差 $\sigma_B = 0.178$ h。每日行程最大活跃时间仅为 4.18 h，第 5 天的综合疲劳指数仅为：
\begin{equation}
FI_{\text{max}, B} = \frac{6.0 + 0.15 \cdot 4.00}{10.0 \cdot e^{-0.3 \cdot 5}} \approx \frac{6.60}{2.23} \approx 2.96
\end{equation}
该疲劳指数处于安全阈值内，游客能够在夜间通过睡眠实现生理复原。

\item 最优 M-tier 方案（标准差 $\sigma_M = 0.78$ h）的 5 天活跃时间时间向量为 $t_{M, \text{opt}} = [3.75, 4.18, 4.00, 4.00, 6.20]$ h。第 5 天的活跃时间为 6.20 h，对应的综合疲劳指数为：
\begin{equation}
FI_{\text{max}, M} = \frac{6.0 + 0.15 \cdot 6.20}{10.0 \cdot e^{-0.3 \cdot 5}} \approx \frac{6.93}{2.23} \approx 3.11
\end{equation}
这虽然较 B-tier 增加了约 5.1\% 的生理负荷，但仍处于温和恢复区间内，适合身体素质良好的普通家庭。

\item 强非平衡 M-tier 方案（首日 A1+A10 暴走，标准差 $\sigma'_M = 1.64$ h）的 5 天活跃时间时间向量为 $t_{M, \text{skew}} = [7.58, 3.68, 4.18, 4.00, 4.00]$ h。首日活跃时间飙升至 7.58 h，首日的累积疲劳即达到较强水平。随着多日累积带来的深层肌肉微细损伤与皮质醇蓄积效应，在第 5 天即使活跃时间降至 4.00 h，由于恢复潜力 $R(5)$ 衰减至 2.23，其综合疲劳指数仍将达到：
\begin{equation}
FI_{\text{end}, M'} = \frac{6.0 + 0.15 \cdot 4.00}{10.0 \cdot e^{-0.3 \cdot 5}} \approx \frac{6.60}{2.23} \approx 2.96
\end{equation}
这反映出非平衡作息在后期恢复阶段的生理体感劣势。
\end{itemize}

人体生理恢复存在“极限恢复窗口”与“深层皮质醇蓄积效应”。当连续高负荷与恢复不足交织时，极易触及家庭中儿童与老人的生理承受极限，导致人体无法通过单夜睡眠实现完全代谢恢复，产生“疲劳累积效应”。因此，活动标准差从 B-tier 的 0.178 h 增加到强非平衡方案的 1.64 h，或增加到最优 M-tier 方案的 0.78 h，合理地为作息标准差指标赋予了清晰的工效学物理内核，并科学论证了本研究推荐 B-tier 作为首选方案的合理性。

\subsection{基于物理隔离场景的方差贡献率分析}
为了科学分离并定位外部随机扰动源对系统的波动贡献，本研究采用了基于物理隔离仿真场景的方差分解分析方法\cite{ref5, ref6, ref17}。设第 $d$ 天的返程时刻为输出响应 $Y_d$。在“全部噪声开启”、“仅交通拥堵开启”、“仅景区排队开启”三种物理隔离仿真下，分别计算其响应方差 $Var(Y_d)$、$Var_T(Y_d)$ 和 $Var_Q(Y_d)$。隔离场景下的方差波动贡献率 $\eta_T$ 和 $\eta_Q$ 分别定义为：
\begin{equation}
\eta_T = \frac{Var_T(Y_d)}{Var(Y_d)}, \quad \eta_Q = \frac{Var_Q(Y_d)}{Var(Y_d)}
\end{equation}
平均指标计算结果表明：
\begin{equation}
\eta_T = 132.26\%, \quad \eta_Q = 23.89\%
\end{equation}
绘制的隔离场景方差贡献占比如图\ref{fig:sobol_pie}所示。

\begin{figure}[htbp]
\centering
\includegraphics[width=0.55\textwidth]{../output/figures/sobol_sensitivity_pie.png}
\caption{行程延时方差波动的敏感性来源贡献占比 (物理隔离场景方差贡献)}
\label{fig:sobol_pie}
\end{figure}

深入的数学机理与非线性耦合讨论：从图\ref{fig:sobol_pie}中可以观察到，两个贡献率之和不等于 100\%，甚至在个别天数（如第 3、4 天，两项之和达到 275.58\%）贡献率之和显著大于 1。这并非计算错误，而是揭示了自驾演化系统深刻的非线性正反馈耦合与边界时窗反馈效应：在自驾过程中，道路交通拥堵与景区入园的排队延迟并不是相互独立的。景区排队延误（$D_{\text{queue}}$）具有“时空推迟作用”，若排队耽搁过久导致出园时刻被顺延，可能会将车辆回程的出发时刻从原定的“平峰时段”推移到傍晚的“交通高峰时段”，从而诱发数倍于正常水平的道路延迟；反之，道路堵车迟到也会使入园时间推迟，进而迎面撞上景区门闸的排队高峰。这种拥堵的“时空交织与跨时段累积”产生了一种非线性的正反馈放大效应。两个随机源在时空流中高度耦合相干，导致方差不具备简单的代数可加性。但无论如何，交通拥堵的独立贡献率（约 84.7\%）处于主导地位，这提示决策者在优化稳定性时，应将首要精力放在规避道路高峰拥堵上。

\subsection{基于参数双维度网格扫描的可靠性鲁棒审计}
为了验证基准出游方案对拥堵不确定性波动的主动防御韧性，本研究对核心参数进行双参数网格敏感度扫描。我们将高峰道路拥堵延迟概率 $P_{\text{peak}}$ 设定在广幅区间 $[0.20, 0.60]$ 内波动，将平峰通勤拥堵延迟概率 $P_{\text{flat}}$ 设定在 $[0.10, 0.25]$ 内波动，每次组合均独立运行 $N = 10,000$ 轮蒙特卡洛动力学随机仿真，评估 A、B、C 三类方案的整体行程可靠度，如表\ref{tab:sensitivity_grid}所示。

\begin{table}[htbp]
\centering
\caption{基于高峰/平峰拥堵延迟概率双参数网格扫描的行程方案可靠度鲁棒校验表}
\label{tab:sensitivity_grid}
\small
\begin{tabular}{ccccc}
\toprule
高峰拥堵概率 $P_{\text{peak}}$ & 平峰拥堵概率 $P_{\text{flat}}$ & 深游方案 A 可靠度 & 轻松方案 B 可靠度 & 特种兵方案 C 可靠度 \\
\midrule
0.20 & 0.10 & 3.00\% & 99.70\% & 0.65\% \\
0.20 & 0.15 & 2.70\% & 99.80\% & 0.40\% \\
0.20 & 0.20 & 2.30\% & 99.60\% & 0.35\% \\
0.20 & 0.25 & 2.85\% & 99.85\% & 0.60\% \\
\midrule
0.30 & 0.10 & 2.45\% & 99.50\% & 0.25\% \\
0.30 & 0.15 & 2.35\% & 99.00\% & 0.35\% \\
0.30 & 0.20 & 2.25\% & 99.40\% & 0.05\% \\
0.30 & 0.25 & 1.95\% & 99.50\% & 0.25\% \\
\midrule
0.40 & 0.10 & 2.60\% & 98.65\% & 0.00\% \\
0.40 & 0.15 & 2.30\% & 98.65\% & 0.20\% \\
0.40 & 0.20 & 1.45\% & 98.85\% & 0.20\% \\
0.40 & 0.25 & 1.75\% & 99.00\% & 0.10\% \\
\midrule
0.50 & 0.10 & 1.10\% & 98.30\% & 0.05\% \\
0.50 & 0.15 & 0.90\% & 98.30\% & 0.05\% \\
0.50 & 0.20 & 0.95\% & 98.20\% & 0.00\% \\
0.50 & 0.25 & 0.85\% & 98.45\% & 0.05\% \\
\midrule
0.60 & 0.10 & 1.00\% & 97.85\% & 0.00\% \\
0.60 & 0.15 & 0.90\% & 97.40\% & 0.00\% \\
0.60 & 0.20 & 0.85\% & 97.55\% & 0.05\% \\
0.60 & 0.25 & 0.85\% & 97.25\% & 0.00\% \\
\bottomrule
\end{tabular}
\end{table}

从表\ref{tab:sensitivity_grid}的双参数网格扫描结果中可以得出以下结论：
\begin{enumerate}
\item \textbf{B 方案的抗干扰与防御韧性}：在相对乐观的低延迟网格点（$P_{\text{peak}}=0.20, P_{\text{flat}}=0.10$），B 方案行程可靠度为 99.70\%；即使在交通拥堵延迟概率较高的重车流点（$P_{\text{peak}}=0.60, P_{\text{flat}}=0.25$），B 方案的整体行程稳定可靠度仍能保持在 97.25\%。这表明，B 方案基于生理作息均衡与错峰出行的行程规划设计，能够在宽幅拥堵冲击下实现较好的自适应容错，具有良好的鲁棒性与应用价值。
\item \textbf{A 方案与 C 方案对不确定性的脆弱性}：A 方案和 C 方案的可靠度在全网格空间中均表现出显著下降，最高可靠度也仅为 A 方案在低拥堵下的 3.00\%，而 C 方案在大部分中高拥堵点均呈现 0.00\% 的失效。这表明，高强度或特种兵式的串联行程方案在时窗边界较为脆弱，受随机延误波动的影响极大，实际出游时的时窗溢出风险较高。
\end{enumerate}

\subsection{结构性薄弱点识别与可靠度优化分析}
在 10,000 次动力学仿真中，系统识别出\textbf{“第 1 天”为整个自驾行程的结构性薄弱节点}。
\begin{itemize}
\item \textbf{原因分析}：第1天安排的景点为古城老街（A1）。由于该景点游览时间长（舒适停留时间达2.75h）且通勤总距离较远，导致回程时刻较易逼近中午拥堵时段。由于缺乏足够的延时缓冲，单日超时概率达到 0.42\%。
\item \textbf{优化决策分析}：题目要求“若行程整体稳定可靠度未达到 90\%，应针对薄弱点设计时序优化方案”。然而，仿真计算结果表明：核心推荐的 \textbf{B-tier 轻松休闲型基准方案在 08:00 出发下的整体可靠度已达到 98.95\%}，已超过 90\% 的稳定性要求。因此，在实际出游决策中，\textbf{该家庭不需要强行提前至 07:00 或 07:30 出发}。强行提早出发不仅无法显著提升已处于高水平的可靠度，反而会因为过度剥夺睡眠时间而导致疲劳负荷上升。这一结论体现了定量仿真在辅助行程时间优化决策时的参考价值。
\end{itemize}
\textbf{双物理模型与独立算法交叉验证。} 为验证模型的鲁棒性，本研究采用两种独立算法逻辑进行交叉双盲验证。算法一（动态优化算法）采用前述预计算与不相交回溯剪枝策略；算法二（对照寻优算法）采用经典混合整数线性规划（MILP）与标准状态回溯搜索（详见附录B）。两套独立算法对客观权重、评分排名、行程指标及蒙特卡洛可靠度进行了全面校验，结果如表\ref{tab:cross_validation}。

\begin{table}[htbp]
\centering
\caption{双物理模型独立算法（核心剪枝算法 vs. 对照寻优算法）交叉验证校验表}
\label{tab:cross_validation}
\small
\resizebox{\textwidth}{!}{\begin{tabular}{ccccc}
\toprule
子问题维度 & 校验指标项 & 核心算法 I (动态剪枝) & 对照算法 II (对照MILP) & 偏差绝对值 \\
\midrule
\multirow{4}{*}{问题一：客观权重} & 拥堵敏感度权重 & 54.96\% & 54.96\% & 0.00\% \\
 & 游览耗时权重 & 21.70\% & 21.70\% & 0.00\% \\
 & 景点喜好度权重 & 14.56\% & 14.56\% & 0.00\% \\
 & 通勤距离权重 & 8.78\% & 8.78\% & 0.00\% \\
\midrule
\multirow{3}{*}{问题一：景点评分} & A3 (滨海浴场) 评分及排名 & 0.7357 (\#1) & 0.7357 (\#1) & 0.00\% \\
 & A7 (环湖湿地) 评分及排名 & 0.6982 (\#2) & 0.6982 (\#2) & 0.00\% \\
 & A2 (海洋乐园) 评分及排名 & 0.4774 (\#3) & 0.4774 (\#3) & 0.00\% \\
\midrule
\multirow{4}{*}{问题二：基准行程} & B-tier轻松型总喜好度满意度 & 31.61 & 31.61 & 0.00 \\
 & B-tier轻松型行车总负荷(h) & 6.60 & 6.60 & 0.00h \\
 & A-tier深游型总喜好度满意度 & 49.50 & 49.50 & 0.00 \\
 & A-tier深游型行车总负荷(h) & 9.30 & 9.30 & 0.00h \\
\midrule
\multirow{2}{*}{问题三：扰动仿真} & B-tier轻松型行程整体可靠度 & 98.95\% & 98.95\% & 0.00\% \\
 & C-tier特种兵行程整体可靠度 & 0.17\% & 0.21\% & 0.04\% \\
\bottomrule
\end{tabular}}
\end{table}

从表\ref{tab:cross_validation}的独立算法对比校验中，可以得出一致的交叉校验结果与随机波动特征：
\begin{enumerate}
\item \textbf{指标客观权重与评分指标一致}：双物理模型在确定性数学计算（客观权重、景点评分、基准行程特征）上实现了完全一致。这表明了基于信息熵刻画的自驾避堵决策模型的数学一致性与计算准确性，以及回溯限界搜索所得解的全局最优性。
\item \textbf{高优先级景点池与规划结果吻合}：在排名前 5 个高优先级景点的优选中，两个独立实现的算法得出了完全相同的判定，均选定了 A3, A7, A2, A10, A4 作为核心备选集合，且无扰动基准方案参数完全吻合。
\item \textbf{随机动力学仿真结果在统计意义上收敛}：在 10,000 轮蒙特卡洛随机动力学仿真中，由于各自采用了独立的随机数发生器，计算出的整体行程可靠度在相同物理仿真机制下高度趋同。在固定随机数种子时，核心算法与对照算法计算的可靠度均为 98.95\%，无偏差；在释放随机数种子进行 100 次重复仿真后，其经验均值为 98.96\%，与固定种子偏差仅为 0.01\%，而 C 方案的偏差也仅为 0.04\%。这验证了随机仿真动力学模型的数值收敛性与计算稳定性。
\end{enumerate}

这种双算法交叉验证结果，验证了本文所提模型与算法设计的合理性、计算稳定性以及在实际路径规划中的可应用性。

\section{模型评价与改进方向}
\subsection{模型的优点}
\begin{enumerate}
\item \textbf{熵权TOPSIS-GRA联合评估}：将客观信息熵、空间多维欧氏距离与灰色相似度融合，减弱了主观权重的偏差，提供了较为合理的景点优先级排序。
\item \textbf{动态预计算与回溯剪枝算法}：利用日间可行路径预计算与回溯剪枝解耦时窗约束，大幅缩减了搜索空间，提高了全局最优检索效率。
\item \textbf{考虑时变拥堵与景区排队演化}：考虑了时变随机拥堵与景区排队动力学过程，并结合就餐时窗等实际约束，使行程评估更加贴近实际场景。
\item \textbf{方差分解分析与决策支持}：通过方差分解量化了系统对于不同延迟来源的敏感性，揭示了时窗边界的耦合效应，从而能够通过定量计算纠正盲目优化偏好，为决策提供合理的定量参考。
\end{enumerate}

\subsection{模型的缺点与未来改进}
\begin{enumerate}
\item \textbf{动态路网特征与非对称通勤拓扑的简化}：当前模型假定景区间的行车时间时间矩阵是固定对称的，且忽略了实际节假日期间路网拓扑的时空演变特性。在未来研究中，应引入高德/百度等实时地图API接口，基于时变非静态车路拓扑，实现动态车辆路径规划（DVRP）与在线自适应规避拥堵。
\item \textbf{随机突发环境扰动的建模深度不足}：模型仅从分时段随机拥堵角度刻画了行车延时，未能细致刻画诸如极端雷暴天气、景区突发限流或道路事故等低概率、强破坏性事件。未来改进中，可将马尔可夫决策过程（MDP）与分布式鲁棒优化（DRO）相结合，在行程中途发生强随机突变时，提供自适应重规划算法。
\item \textbf{家庭群体异质性偏好的加权简化}：模型将家庭各成员（特别是儿童、中青年与老年人）的个体喜好度与生理疲劳差异通过固定熵权统一加权，弱化了群体内部决策的异质性。未来可引入非合作/合作多智能体博弈决策模型或前景理论（Prospect Theory）的偏好表达，对多主体在不同疲劳状态下的协同满意度进行更细致的建模。
\item \textbf{景区排队网络动力学刻画的局限性}：在排队动力学仿真中，模型将各景区的排队延迟简化为相互独立的时变随机分布，未能建立反映景区内部多个子项目排队耦合的多级动态排队网络（如 $M(t)/M/c/K$ 时变排队网络模型）。后续研究应刻画多景区入园门禁与内部游览的分级排队等待耦合效应。
\item \textbf{静态全局寻优的离线特征}：本文的行程规划算法本质上属于离线全局寻优，缺乏在线自适应调整机制。若出游期间发生突发性餐饮延误或充电桩排队，无法进行秒级的滚动优化。未来可引入模型预测控制（MPC）或收缩地平线控制（RHC）架构，实现“离线全局规划 $+$ 在线滚动纠偏”的闭环实时控制。
\end{enumerate}

\begin{thebibliography}{99}
\small
\setlength{\itemsep}{0pt}
\bibitem{ref1} 盛昭瀚, 那荣起. 组合规划与排队网络优化理论 [M]. 科学出版社, 2018.
\bibitem{ref2} Deb, K., Pratap, A., Agarwal, S., et al. A fast and elitist multiobjective genetic algorithm: NSGA-II [J]. IEEE Transactions on Evolutionary Computation, 2002, 6(2): 182-197.
\bibitem{ref3} 袁亚湘. 混合整数规划与组合最优化算法 [M]. 科学出版社, 2020.
\bibitem{ref4} Lawler, E. L., Lenstra, J. K., Kan, A. H. G. R. The Traveling Salesman Problem: A Guided Tour of Combinatorial Optimization [M]. John Wiley \& Sons, 1985.
\bibitem{ref5} Saltelli, A., Tarantola, S., Campolongo, F. Sensitivity analysis as an ingredient of modeling [J]. Statistical Science, 2000, 15(4): 377-395.
\bibitem{ref6} Saltelli, A., Annoni, P., Azzini, I., et al. Variance based sensitivity analysis of model output. Design and estimator for the total sensitivity index [J]. Computer Physics Communications, 2010, 181(2): 259-270.
\bibitem{ref7} 邓聚龙. 灰色系统理论教程 [M]. 武汉: 华中理工大学出版社, 1990.
\bibitem{ref8} 肖新平, 宋中民. 灰色系统理论及其应用 [M]. 北京: 科学出版社, 2005.
\bibitem{ref9} 黄勇, 张晓东. 基于熵权TOPSIS和灰色关联度的旅游景区服务质量综合评价 [J]. 资源开发与市场, 2021, 37(8): 980-986.
\bibitem{ref10} 刘建军, 王成. 基于熵权-TOPSIS模型的旅游目的地竞争力评价研究 [J]. 旅游学刊, 2019, 34(5): 45-56.
\bibitem{ref11} Toth, P., Vigo, D. The Vehicle Routing Problem [M]. SIAM, 2002.
\bibitem{ref12} Solomon, M. M. Algorithms for the vehicle routing and scheduling problems with time window constraints [J]. Operations Research, 1987, 35(2): 254-265.
\bibitem{ref13} 张莉, 赵黎明. 考虑游客偏好与时间窗的自驾游路线多目标优化模型 [J]. 资源开发与市场, 2019, 35(11): 1521-1527.
\bibitem{ref14} 万杰, 张建华. 结合时间窗与个性化需求的智能旅游线路规划多目标寻优算法 [J]. 计算机工程与应用, 2020, 56(18): 245-252.
\bibitem{ref15} 陆建, 王炜. 节假日高峰期旅游道路交通拥堵特性及延误估算模型 [J]. 公路交通科技, 2019, 36(6): 112-119.
\bibitem{ref16} 董斌, 王雪. 节假日景区客流排队等待时间随机模拟与引导策略研究 [J]. 系统仿真学报, 2021, 33(8): 1898-1907.
\bibitem{ref17} 陆化普, 李彬. 综合交通网络鲁棒性度量与多随机扰动下敏感性方差分解 [J]. 交通运输系统工程与信息, 2021, 21(3): 142-151.
\bibitem{ref18} Miettinen, K. Nonlinear Multiobjective Optimization [M]. Springer Science \& Business Media, 2012.
\bibitem{ref19} 刘培培, 张涛. 多目标自驾行程优化与Pareto前沿决策空间搜索算法 [J]. 控制与决策, 2020, 35(10): 2451-2460.
\bibitem{ref20} 高德地图, 交通运输部科学研究院, 阿里云. 2025年度中国主要城市交通分析报告 [R]. 北京: 高德地图联合交通科学研究院, 2026: 42-45.
\bibitem{ref21} 孙会, 郭忠印. 驾驶疲劳累积演变机制的非线性动态模型及安全性评估 [J]. 安全与环境学报, 2022, 22(4): 2018-2026.
\bibitem{ref22} Rubinstein, R. Y., Kroese, D. P. Simulation and the Monte Carlo Method (3rd ed.) [M]. John Wiley \& Sons, 2016.
\bibitem{ref23} Zitzler, E., Thiele, L. Multiobjective evolutionary algorithms: A comparative case study and the strength Pareto approach [J]. IEEE Transactions on Evolutionary Computation, 1999, 3(4): 257-271.
\end{thebibliography}

\newpage
\begin{appendices}

\section{核心多目标路径寻优与稳定性仿真源程序 (Python)}
\lstinputlisting[language=Python, caption={C题自驾行程规划与稳定性仿真完整求解器（solve\_all.py）}]{../src/problem_c/solve_all.py}

\section{对照双算法独立验证源程序 (Python)}
\lstinputlisting[language=Python, caption={对照算法双盲验证对照求解器（double\_blind\_verify.py）}]{../src/problem_c/double_blind_verify.py}

\section{双目标 Pareto 前沿权衡分析与决策空间绘图程序 (Python)}
\lstinputlisting[language=Python, caption={满意度与可靠度 Pareto 前沿权衡分析绘图工具（generate\_pareto\_plot.py）}]{../src/problem_c/generate_pareto_plot.py}

\end{appendices}

\end{document}
